\documentclass[10pt]{article}
\usepackage[T1]{fontenc}
\usepackage[utf8]{inputenc}
\usepackage[french]{babel}
\usepackage{amsmath,amssymb,mathrsfs,mathpazo}
\usepackage{geometry}
\geometry{a4paper,margin=2cm}
\usepackage{graphicx,array,tabularx,tabularray,colortbl,multirow,multicol,booktabs,enumitem,pifont,ulem,fancybox,diagbox,cancel}
\usepackage{pstricks,pstricks-add,pst-plot,pst-tree,pst-node,pst-text,pst-func,pst-math,pst-bspline,pst-3dplot,pst-eucl,pst-solides3d}
\let\pscalculate\undefined
\usepackage{pst-bezier}
\usepackage{tikz,pgfplots,tkz-tab}
\usetikzlibrary{arrows,arrows.meta,patterns,calc,positioning,trees,decorations.markings,shapes.geometric,shapes.misc,shapes.arrows,intersections,decorations.pathmorphing,decorations.pathreplacing,backgrounds,shadows,babel}
\pgfplotsset{compat=1.18}
\usepgfplotslibrary{fillbetween}
\usepackage[np]{numprint}
\usepackage{listings,makecell,pas-tableur,twoopt}
\let\barre\undefined
\usepackage[pstricks]{tablvar}
\definecolor{darkgreen}{RGB}{0,100,0}
\providecommand{\red}{\color{red}}
\providecommand{\blue}{\color{blue}}
\providecommand{\green}{\color{green}}
\providecommand{\black}{\color{black}}
\providecommand{\white}{\color{white}}
\providecommand{\gray}{\color{gray}}
\providecommand{\cyan}{\color{cyan}}
\providecommand{\magenta}{\color{magenta}}
\providecommand{\yellow}{\color{yellow}}
\providecommand{\violet}{\color{violet}}
\providecommand{\orange}{\color{orange}}
\providecommand{\euro}{\texteuro}
\providecommand{\EUR}{\texteuro}
\providecommand{\ieme}{\textsuperscript{e}}
\providecommand{\degres}{\ensuremath{^\circ}}
\providecommand{\degree}{\ensuremath{^\circ}}
\providecommand{\vv}[1]{\overrightarrow{#1}}
\providecommand{\Coord}[1]{#1}
\providecommand{\decosix}{}
\providecommand{\starredbullet}{\ensuremath{\bullet}}
\pagestyle{empty}
\setlength{\parindent}{0pt}
\usepackage{adjustbox}\DeclareUnicodeCharacter{2236}{:}\DeclareUnicodeCharacter{2190}{\ensuremath{\leftarrow}}
\providecommand{\R}{\mathbb{R}}
\providecommand{\N}{\mathbb{N}}
\providecommand{\Z}{\mathbb{Z}}
\providecommand{\Q}{\mathbb{Q}}
\providecommand{\D}{\mathbb{D}}
\providecommand{\C}{\mathbb{C}}
\providecommand{\e}{{\mathrm{e}}}
\providecommand{\vect}[1]{\overrightarrow{#1}}
\providecommand{\vectt}[1]{\overrightarrow{\mathrm{#1}}}
\providecommand{\vvt}[1]{\overrightarrow{\mathrm{#1}}}
\providecommand{\np}[2][]{{#2~#1}}
\providecommand{\vv}[1]{{\overrightarrow{#1}}}
\renewcommand{\d}{{\mathrm{d}}}
\providecommand{\degres}{{^\circ}}
\providecommand{\degre}{{^\circ}}
\providecommand{\blue}{\color{blue}}
\providecommand{\red}{\color{red}}
\providecommand{\enskip}{\quad}
\providecommand{\llbracket}{\lbrack\!\lbrack}
\providecommand{\rrbracket}{\rbrack\!\rbrack}
\providecommand{\euro}{{\text{€}}}
\providecommand{\no}{{\text{n°}}}
\providecommand{\interval}[2]{{[#1;#2]}}
\providecommand{\intervFO}[2]{{[#1;#2[}}
\providecommand{\dotfill}{\ldots}

\providecommand{\ds}{\displaystyle}
\providecommand{\textding}[1]{#1}
\providecommand{\Oij}{\(\mathrm{O};\vec{\imath},\vec{\jmath}\)}
\providecommand{\Oijk}{\(\mathrm{O};\vec{\imath},\vec{\jmath},\vec{k}\)}
\providecommand{\Ouv}{\(\mathrm{O};\vec{u},\vec{v}\)}
\providecommand{\intervFF}[2]{\ensuremath{[#1;#2]}}
\renewcommand{\euro}{\ifmmode\text{\texteuro}\else\texteuro\fi}
\DeclareUnicodeCharacter{00B0}{\ifmmode^{\circ}\else\textdegree\fi}
\DeclareUnicodeCharacter{20AC}{\ifmmode\text{\texteuro}\else\texteuro\fi}
% Énoncés : Éducation nationale. Transcription et corrigés : APMEP, Denis Vergès et contributeurs. Découpage et classement : CoopMaths, Sébastien Lozano, Éric Elter et contributeurs. Adaptation de la présentation : Maths Explorer.
% CC BY-SA : https://coopmaths.fr/www/about/
% https://coopmaths.fr/alea/static/e3c/2020/tex/e3c_2020_00_sujet49_4_cor.tex
\begin{document}
\noindent\textit{APMEP / CoopMaths. CC BY-SA. Présentation adaptée par Maths Explorer.}\par\medskip
\( f(x) = (ax + b)\e^{-0{,}1x} \).

\subsection*{1.}

On calcule \( f(0) = (0 + b)\e^{-0{,}1 \times 0} = b\e^0 = b = 5 \) puisque \( A(0\,;\,5) \in \mathcal{C}_f \).

\subsection*{2.}

\paragraph{a.} Pour l'équation réduite de \( (AB) \), l'ordonnée à l'origine est 5 et le coefficient directeur de la droite est égal à :
\[
\dfrac{19 - 5}{4 - 0} = \dfrac{14}{4} = 3{,}5.
\]
L'équation réduite de la droite \( (AB) \) est donc \( y = 3{,}5x + 5 \).

\paragraph{b.} On a \( f(x) = (ax + 5) \e^{-0{,}1x} \), \( f \) est un produit de fonctions dérivables sur \( \mathbb{R} \), donc en dérivant comme un produit \( ( (uv)' = u'v + uv' ) \) :
\begin{align*}
f'(x) &= a\e^{-0{,}1x} - 0{,}1(ax + 5)\e^{-0{,}1x} \\
&= \e^{-0{,}1x} [a - 0{,}1(ax + 5)] \\
&= \e^{-0{,}1x}(a - 0{,}1ax - 0{,}5).
\end{align*}
Or on sait que le coefficient directeur de la tangente à la courbe au point d'abscisse 0 est le nombre dérivé \( f'(0) \). Donc :
\[
f'(0) = \e^0(a - 0 - 0{,}5) = a - 0{,}5 = 3{,}5, \text{ d'où on déduit } a = 3{,}5 + 0{,}5 = 4.
\]
On a donc sur \( \mathbb{R} \), \( f(x) = (4x + 5)\e^{-0{,}1x} \).

\subsection*{3.}

\paragraph{a.} Avec l'expression trouvée pour \( f'(x) \) ci-dessus et en remplaçant \( a \) par 4, on obtient :
\[
f'(x) = \e^{-0{,}1x}(4 - 0{,}1 \times 4x - 0{,}5) = (-0{,}4x + 3{,}5)\e^{-0{,}1x}.
\]

\paragraph{b.} On sait que, quel que soit \( x \in \mathbb{R} \), \( \e^{-0{,}1x} > 0 \), donc le signe de \( f'(x) \) est celui de \( -0{,}4x + 3{,}5 \).
\begin{align*}
&-0{,}4x + 3{,}5 > 0  \\
\iff &3{,}5 > 0{,}4x \\
\iff &\dfrac{3{,}5}{0{,}4} > x \\
\iff &x < 8{,}75
\end{align*}
\( f \) est croissante sur \( ] -\infty \,; 8{,}75[ \), puis décroissante sur \( ]8{,}75 \,; +\infty[ \), avec un maximum en :
\[
f(8{,}75) = (4 \times 8{,}75 + 5)\e^{-0{,}1 \times 8{,}75} \approx 16{,}675.
\]
(ce que confirme la figure)

\end{document}
